% 可直接插入主文档的组合数 LaTeX 片段
% 导言区需要加载：ctex、amsmath、amssymb、tcolorbox、minted

\section{组合数}

\subsection{常见公式和经典问题}

记
\[
    \binom{n}{m}=C_n^m=\frac{n!}{m!(n-m)!},\qquad
    A_n^m=\frac{n!}{(n-m)!}=\binom{n}{m}m!.
\]
当 $m<0$ 或 $m>n$ 时，约定 $\binom{n}{m}=0$。

\begin{itemize}
    \item \textbf{组合数：}从 $n$ 个不同物品中选出 $m$ 个，不计顺序，方案数为
    $\binom{n}{m}$。
    \item \textbf{排列数：}从 $n$ 个不同物品中选出 $m$ 个并排列，方案数为
    $A_n^m=\binom{n}{m}m!$。
    \item \textbf{二项式定理：}
    \[
        (x+y)^n=\sum_{k=0}^{n}\binom{n}{k}x^ky^{n-k}.
    \]
    \item \textbf{多项式定理：}
    \[
        (x_1+x_2+\cdots+x_m)^n
        =\sum_{k_1+\cdots+k_m=n}
        \frac{n!}{k_1!k_2!\cdots k_m!}
        x_1^{k_1}x_2^{k_2}\cdots x_m^{k_m}.
    \]
    \item 从 $n$ 个物品中可重复取得 $k$ 个，每次有顺序，方案数为 $n^k$。
    \item 从 $n$ 个物品中不可重复取 $k$ 个做排列，方案数为
    $A_n^k=\binom{n}{k}k!$。
    \item 从 $n$ 个物品中选 $k$ 个做圆排列，若只考虑旋转同构，方案数为
    $\binom{n}{k}(k-1)!=\dfrac{n!}{k(n-k)!}$。
    若翻转也视为相同，通常在 $k>2$ 时再除以 $2$。
    \item $n$ 个物品中，第 $i$ 种有 $k_i$ 个且 $\sum_{i=1}^m k_i=n$，
    全部排列数为
    \[
        \frac{n!}{k_1!k_2!\cdots k_m!}.
    \]
    \item 从 $n$ 种物品中可重复选 $k$ 个做组合，方案数为
    \[
        \binom{n+k-1}{k}=\binom{n+k-1}{n-1}.
    \]
    \item 从 $n$ 个位置中选 $k$ 个互不相邻的位置，方案数为
    $\binom{n-k+1}{k}$。
    \item 从 $n$ 个位置中选 $k$ 个，要求首尾也不能同时选，方案数为
    \[\binom{n-k+1}{k}-\binom{n-k-1}{k-2}.\]
\end{itemize}

\subsection{基本性质}

\begin{itemize}
    \item \textbf{对称性：}
    \[
        \binom{n}{k}=\binom{n}{n-k}.
    \]
    \item \textbf{Pascal 递推：}
    \[
        \binom{n}{k}=\binom{n-1}{k-1}+\binom{n-1}{k}.
    \]
    \item \textbf{边界：} $\binom{n}{0}=\binom{n}{n}=1$。
    \item \textbf{相邻项比值：}
    \[
        \frac{\binom{n}{k+1}}{\binom{n}{k}}=\frac{n-k}{k+1}.
    \]
    \item \textbf{组合数递推：}
    \[
        k\binom{n}{k}=n\binom{n-1}{k-1},\qquad
        (n-k)\binom{n}{k}=n\binom{n-1}{k}.
    \]
    \item \textbf{先选后选：}
    \[
        \binom{n}{k}\binom{k}{r}
        =\binom{n}{r}\binom{n-r}{k-r}.
    \]
    \item \textbf{下降幂：}
    \[
        k^{\underline r}=k(k-1)\cdots(k-r+1)=r!\binom{k}{r}.
    \]
\end{itemize}

\subsection{插板法}

非负整数解的个数为
\[
    x_1+x_2+\cdots+x_n=k,\quad x_i\ge0
    \quad\Longrightarrow\quad
    \binom{n+k-1}{n-1}.
\]
正整数解的个数为
\[
    x_1+x_2+\cdots+x_n=k,\quad x_i\ge1
    \quad\Longrightarrow\quad
    \binom{k-1}{n-1}.
\]
若 $x_i\ge a_i$，先令 $y_i=x_i-a_i$ 再使用插板法。

带上界时常用容斥：若 $0\le x_i\le b_i$，则
\[
\#\{x_1+\cdots+x_n=k\}
 =\sum_{S\subseteq\{1,\dots,n\}}(-1)^{|S|}
 \binom{k-\sum_{i\in S}(b_i+1)+n-1}{n-1}.
\]

\subsection{经典恒等式}

\begin{itemize}
    \item \textbf{二项式求和：}
    \[
        \sum_{k=0}^{n}\binom{n}{k}=2^n,
        \qquad
        \sum_{k=0}^{n}(-1)^k\binom{n}{k}=0\quad(n>0).
    \]
    \item \textbf{带权二项式和：}
    \[
        \sum_{k=0}^{n}x^k\binom{n}{k}=(1+x)^n.
    \]
    特别地，$\sum_{k=0}^n2^k\binom{n}{k}=3^n$。
    \item \textbf{一次加权和：}
    \[
        \sum_{k=0}^{n}k\binom{n}{k}=n2^{n-1}.
    \]
    \item \textbf{二次加权和：}
    \[
        \sum_{k=0}^{n}k(k-1)\binom{n}{k}=n(n-1)2^{n-2},
    \]
    \[
        \sum_{k=0}^{n}k^2\binom{n}{k}=n(n+1)2^{n-2}.
    \]
    \item \textbf{一般下降幂加权和：}
    \[
        \sum_{k=0}^{n}k^{\underline r}\binom{n}{k}
        =n^{\underline r}2^{n-r}.
    \]
    \item \textbf{有限差分：}
    \[
        \sum_{k=0}^{n}(-1)^k\binom{n}{k}k^m=0\quad(m<n),
    \]
    \[
        \sum_{k=0}^{n}(-1)^k\binom{n}{k}k^n=(-1)^n n!.
    \]
    \item \textbf{平方和：}
    \[
        \sum_{k=0}^{n}\binom{n}{k}^2=\binom{2n}{n}.
    \]
    \item \textbf{相邻平方和：}
    \[
        \sum_{k=0}^{n}\binom{n}{k}\binom{n}{k+1}
        =\binom{2n}{n-1}.
    \]
    \item \textbf{冰球杆恒等式：}
    \[
        \sum_{i=r}^{n}\binom{i}{r}=\binom{n+1}{r+1}.
    \]
\end{itemize}

\subsection{范特蒙德卷积及其变形}

\textbf{范特蒙德卷积：}
\[
    \sum_{i=0}^{k}\binom{r}{i}\binom{s}{k-i}
    =\binom{r+s}{k}.
\]
它表示从 $r+s$ 个物品中选 $k$ 个，按前 $r$ 个中选了多少分类。

常用变形包括
\[
    \sum_{i=0}^{n}\binom{r}{i}\binom{s}{n-i}=\binom{r+s}{n},
\]
\[
    \sum_i(-1)^i\binom{r}{i}\binom{s}{k-i}
    =[x^k](1-x)^r(1+x)^s.
\]
特别地，利用 $\binom{s}{i}=\binom{s}{s-i}$ 可得
\[
    \sum_{i=0}^{\min(r,s)}\binom{r}{i}\binom{s}{i}
    =\binom{r+s}{r}.
\]

特别常用的形式为
\[
    \sum_{i=0}^{n}\binom{n}{i}\binom{m}{k-i}
    =\binom{n+m}{k},
\]
\[
    \sum_{i=0}^{n}\binom{n}{i}^2=\binom{2n}{n},
    \qquad
    \sum_{i=0}^{n}\binom{n}{i}\binom{n}{i+1}=\binom{2n}{n-1}.
\]
还有
\[
    \sum_{i=0}^{n}\binom{i}{r}\binom{n-i}{s}
    =\binom{n+1}{r+s+1}.
\]

\subsection{二项式反演}

若
\[
    f(n)=\sum_{k=0}^{n}\binom{n}{k}g(k),
\]
则
\[
    g(n)=\sum_{k=0}^{n}(-1)^{n-k}\binom{n}{k}f(k).
\]

等价地，若
\[
    f(n)=\sum_{k=0}^{n}(-1)^k\binom{n}{k}g(k),
\]
则
\[
    g(n)=\sum_{k=0}^{n}(-1)^k\binom{n}{k}f(k).
\]

\subsection{容斥原理}

有限集合的并集大小为
\[
    \left|\bigcup_{i=1}^{n}A_i\right|
    =\sum_{\emptyset\ne S\subseteq\{1,\dots,n\}}
    (-1)^{|S|+1}\left|\bigcap_{i\in S}A_i\right|.
\]

也可以写成按交集大小分组的形式：
\[
    \left|\bigcup_iA_i\right|
    =\sum_i|A_i|-\sum_{i<j}|A_i\cap A_j|
    +\sum_{i<j<k}|A_i\cap A_j\cap A_k|-\cdots.
\]

全集 $U$ 中同时满足所有条件 $S_i$ 的元素个数为
\[
    \left|\bigcap_{i=1}^{n}S_i\right|
    =|U|-\left|\bigcup_{i=1}^{n}\overline{S_i}\right|.
\]

\subsection{Catalan 数}

第 $n$ 个 Catalan 数为
\[
    \mathrm{Cat}_n
    =\frac{1}{n+1}\binom{2n}{n}
    =\binom{2n}{n}-\binom{2n}{n+1}.
\]
常见模型包括：合法括号序列、$n+1$ 个叶子的不同二叉树、凸多边形三角剖分、
从左上到右下且不越过对角线的路径等。

递推式为
\[
    \mathrm{Cat}_0=1,\qquad
    \mathrm{Cat}_{n+1}=\sum_{i=0}^{n}\mathrm{Cat}_i\mathrm{Cat}_{n-i}.
\]

\subsection{错排数}

$n$ 个元素的错排数记为 $D_n$，即每个元素都不在原位置：
\[
    D_n=n!\sum_{i=0}^{n}\frac{(-1)^i}{i!}.
\]
常用递推为
\[
    D_n=(n-1)(D_{n-1}+D_{n-2}),
    \qquad D_0=1,\ D_1=0.
\]

\subsection{其他常用模型}

\begin{itemize}
    \item 长度为 $n$ 的 01 串中恰有 $k$ 个 $1$：$\binom{n}{k}$。
    \item 长度为 $n$ 的 01 串中恰有 $k$ 个连续段：
    $2\binom{n-1}{k-1}$。
    \item 从 $1,2,\dots,n$ 中选 $k$ 个不相邻数：$\binom{n-k+1}{k}$。
    \item 斐波那契数的组合表示：
    \[
        F_{n+1}=\sum_{k\ge0}\binom{n-k}{k}.
    \]
    \item $n$ 个不同元素的圆排列：$(n-1)!$；若还要区分翻转，需根据题意处理镜像。
\end{itemize}

\subsection{询问排列数、组合数}

当模数 $\mathrm{mod}=998244353$ 为质数且 $n<\mathrm{mod}$ 时，预处理阶乘和
阶乘逆元即可在 $O(1)$ 时间回答组合数与排列数询问。

\begin{minted}{cpp}
const int N=1000010;
const int mod=998244353;
int fac[N],ifac[N];

int qpow(int a,int b)
{
    int res=1;
    while(b)
    {
        if(b&1)res=(ll)res*a%mod;
        a=(ll)a*a%mod;
        b>>=1;
    }
    return res;
}

void init()
{
    fac[0]=ifac[0]=1;
    for(int i=1;i<N;i++)fac[i]=(ll)fac[i-1]*i%mod;
    ifac[N-1]=qpow(fac[N-1],mod-2);
    for(int i=N-2;i>=0;i--)
        ifac[i]=(ll)ifac[i+1]*(i+1)%mod;
}

int C(int n,int m)
{
    if(m<0||m>n)return 0;
    return (ll)fac[n]*ifac[m]%mod*ifac[n-m]%mod;
}

int P(int n,int m)
{
    if(m<0||m>n)return 0;
    return (ll)fac[n]*ifac[n-m]%mod;
}
\end{minted}

\subsection{Lucas 定理}

当 $n$ 很大而质数模数 $p$ 较小时，Lucas 定理为
\[
    \binom{n}{m}\bmod p
    =\binom{\lfloor n/p\rfloor}{\lfloor m/p\rfloor}
    \binom{n\bmod p}{m\bmod p}\bmod p.
\]
递归次数为 $O(\log_p n)$，每次组合数在 $[0,p)$ 内计算。
下面的 Lucas 模板要求先把 $fac,ifac$ 预处理到 $p-1$，因此适合质数模数较小的情形。

\begin{minted}{cpp}
int Csmall(int n,int m)
{
    if(m<0||m>n)return 0;
    return (ll)fac[n]*ifac[m]%mod*ifac[n-m]%mod;
}

int Lucas(long long n,long long m)
{
    if(m<0||m>n)return 0;
    if(!m)return 1;
    return (ll)Csmall(n%mod,m%mod)
        *Lucas(n/mod,m/mod)%mod;
}
\end{minted}

\subsection{Pascal 三角形}

当模数不适合使用逆元，或 $n$ 较小时，可以直接使用 Pascal 递推：
\[
    \binom{i}{0}=\binom{i}{i}=1,\qquad
    \binom{i}{j}=\binom{i-1}{j-1}+\binom{i-1}{j}.
\]
时间复杂度为 $O(n^2)$，空间可以滚动到 $O(n)$。

\begin{minted}{cpp}
for(int i=0;i<=n;i++)
{
    C[i][0]=C[i][i]=1;
    for(int j=1;j<i;j++)
        C[i][j]=(C[i-1][j-1]+C[i-1][j])%mod;
}
\end{minted}

\subsection{常用结论速查}

\begin{center}
\begin{tabular}{c|c}
\hline
问题 & 答案 \\
\hline
从 $n$ 个不同物品选 $k$ 个 & $\binom{n}{k}$ \\
从 $n$ 个不同物品选 $k$ 个排列 & $\dfrac{n!}{(n-k)!}$ \\
非负解 $x_1+\cdots+x_n=k$ & $\binom{n+k-1}{n-1}$ \\
正整数解 $x_1+\cdots+x_n=k$ & $\binom{k-1}{n-1}$ \\
$n$ 个位置选 $k$ 个不相邻 & $\binom{n-k+1}{k}$ \\
圆排列 $n$ 个不同物品 & $(n-1)!$ \\
错排 & $n!\sum_{i=0}^{n}(-1)^i/i!$ \\
Catalan & $\dfrac{1}{n+1}\binom{2n}{n}$ \\
\hline
\end{tabular}
\end{center}
